% The blueprint: the mathematics, in LaTeX, with each statement linked to Lean.
%   \lean{Name}   the Lean declaration that formalizes it (checked by checkdecls)
%   \leanok       in a statement: it is formalized; in a proof: the proof is done
%   \uses{label}  what it depends on -- the arrows of the dependency graph

\chapter{Odd numbers add up to squares}

\begin{definition}[Sum of odd numbers]\label{def:oddSum}
  \lean{OddSquares.oddSum}\leanok
  For $n \in \mathbb{N}$, let $S(n) = 1 + 3 + \dots + (2n - 1)$ be the sum of the first $n$ odd
  numbers, with $S(0) = 0$.
\end{definition}

\begin{lemma}\label{lem:oddSum_succ}
  \lean{OddSquares.oddSum_succ}\leanok
  \uses{def:oddSum}
  $S(n + 1) = S(n) + (2n + 1)$.
\end{lemma}
\begin{proof}\leanok
  The next odd number after the first $n$ is $2n + 1$.
\end{proof}

\begin{theorem}\label{thm:oddSum_eq_sq}
  \lean{OddSquares.oddSum_eq_sq}\leanok
  \uses{def:oddSum}
  $S(n) = n^2$ for every $n$.
\end{theorem}
\begin{proof}\leanok
  \uses{lem:oddSum_succ}
  By induction on $n$. $S(0) = 0 = 0^2$, and
  $S(n + 1) = S(n) + 2n + 1 = n^2 + 2n + 1 = (n + 1)^2$.
\end{proof}

\begin{corollary}\label{cor:oddSum_isSquare}
  \lean{OddSquares.oddSum_isSquare}\leanok
  \uses{def:oddSum}
  $S(n)$ is a perfect square.
\end{corollary}
\begin{proof}\leanok
  \uses{thm:oddSum_eq_sq}
  Take $m = n$.
\end{proof}

\chapter{Sums of cubes}

\begin{definition}[Sum of the first numbers]\label{def:sumTo}
  \lean{OddSquares.sumTo}\leanok
  $T(n) = 0 + 1 + \dots + n$.
\end{definition}

\begin{theorem}[Gauss]\label{thm:two_mul_sumTo}
  \lean{OddSquares.two_mul_sumTo}\leanok
  \uses{def:sumTo}
  $2\,T(n) = n(n + 1)$.
\end{theorem}
\begin{proof}\leanok
  By induction on $n$: $2\,T(n + 1) = n(n + 1) + 2(n + 1) = (n + 1)(n + 2)$.
\end{proof}

\begin{definition}[Sum of cubes]\label{def:cubeSum}
  \lean{OddSquares.cubeSum}\leanok
  $C(n) = 0^3 + 1^3 + \dots + n^3$.
\end{definition}

\begin{theorem}[Nicomachus]\label{thm:nicomachus}
  \lean{OddSquares.cubeSum_eq_sq_sumTo}\leanok
  \uses{def:cubeSum, def:sumTo}
  $C(n) = T(n)^2$: the sum of the first cubes is the square of the sum of the first numbers.
\end{theorem}
\begin{proof}
  \uses{thm:two_mul_sumTo}
  By induction on $n$. The step needs $(n + 1)^3 = T(n + 1)^2 - T(n)^2$, which follows from
  Gauss's formula. This proof is not formalized yet: in Lean it is still \texttt{sorry}.
\end{proof}
